\documentclass[11pt,a4paper]{article}
\usepackage[T1]{fontenc}
\usepackage[utf8]{inputenc}
\usepackage[main=italian,provide=*]{babel}
\usepackage{amsmath,amssymb}
\usepackage{geometry}
\geometry{margin=2.3cm,top=2.5cm,bottom=2.7cm}
\usepackage{booktabs}
\usepackage[table,dvipsnames]{xcolor}
\usepackage{tcolorbox}
\tcbuselibrary{skins,breakable}
\usepackage{titlesec}
\usepackage{enumitem}
\usepackage{needspace}
\usepackage{fancyhdr}
\usepackage{graphicx}
\usepackage{hyperref}
\usepackage{longtable}
\hypersetup{colorlinks=true,linkcolor=Sepia,citecolor=Sepia,urlcolor=Sepia}
\newcommand{\Tr}{\operatorname{Tr}}
\newcommand{\ket}[1]{\lvert #1 \rangle}

\definecolor{inkblue}{HTML}{1B3A5C}
\definecolor{lightblue}{HTML}{EAF1F8}
\definecolor{accent}{HTML}{B0562A}
\definecolor{softgray}{HTML}{6B6B6B}
\definecolor{okgreen}{HTML}{2E6B3E}
\definecolor{okgreenbg}{HTML}{EAF5EC}
\definecolor{warnamber}{HTML}{9C6B12}
\definecolor{warnbg}{HTML}{FBF1E0}

\titleformat{\section}
  {\normalfont\Large\bfseries\color{inkblue}}
  {\thesection}{0.6em}{}[{\color{inkblue!40}\titlerule}]
\titlespacing{\section}{0pt}{0.9em}{0.4em}
\titleformat{\subsection}
  {\normalfont\large\bfseries\color{accent}}
  {}{0em}{}
\titlespacing{\subsection}{0pt}{0.6em}{0.2em}

\pagestyle{fancy}
\fancyhf{}
\renewcommand{\headrulewidth}{0.4pt}
\fancyhead[L]{\small\color{softgray}Guida d'uso --- pacchetto riproducibile, dimero (Parte 1)}
\fancyhead[R]{\small\color{softgray}\thepage}
\fancyfoot[C]{\small\color{softgray}Tesi triennale, Samuele Schiavi --- Relatore: Prof.\ Alessandro Chiesa}

\newtcolorbox{keyeq}[1][]{
  colback=lightblue, colframe=inkblue!70, boxrule=0.6pt,
  arc=2pt, left=8pt, right=8pt, top=3pt, bottom=3pt, #1
}
\newtcolorbox{okbox}[1]{
  colback=okgreenbg, colframe=okgreen!60, boxrule=0.6pt,
  arc=2pt, left=7pt, right=7pt, top=2pt, bottom=2pt,
  fonttitle=\bfseries\color{okgreen}, title={#1}
}
\newtcolorbox{warnbox}[1]{
  colback=warnbg, colframe=warnamber!70, boxrule=0.6pt,
  arc=2pt, left=7pt, right=7pt, top=2pt, bottom=2pt,
  fonttitle=\bfseries\color{warnamber}, title={#1}
}

\newcommand{\code}[1]{\texttt{\small #1}}

\title{\vspace{-1.5em}\color{inkblue}Guida d'uso\\[0.2em]
\Large\normalfont\color{softgray}Riprodurre il dimero senza rumore (Parte 1) in locale}
\author{Pacchetto riproducibile --- Tesi triennale, Samuele Schiavi\\ Relatore: Prof.\ Alessandro Chiesa}
\date{}

\begin{document}
\sloppy
\maketitle
\thispagestyle{fancy}
\vspace{-2.0em}

Questo pacchetto contiene tutto il codice necessario per riprodurre, da
zero, la Parte 1 (sistema chiuso, nessun rumore) del lavoro sul dimero:
diagonalizzazione esatta, confronto fra ansatz VQE, evoluzione di
Trotter e correlatori dinamici. E' il pacchetto \emph{fratello} di
\code{pacchetto\_dimero\_rumoroso\_riproducibile} (Parte 2) --- questo
copre solo il sistema chiuso, non il rumore. Ogni modulo e ogni funzione
sono commentati esplicitando sia il loro ruolo all'interno del proprio
file sia il loro ruolo nell'insieme della pipeline.

\tableofcontents

\section{Struttura del pacchetto}

\begin{keyeq}
\centering
\begin{tabular}{@{}lp{9.5cm}@{}}
\code{codice/} & 9 moduli Python (mattoni + confronto ansatz + Trotter + correlatori)\\
\code{dati/} & dati intermedi (\code{.npz}) generati da 01\_, 02\_ e 03\_\\
\code{figure/} & 19 figure finali, in PDF e PNG (15 aggiunte in questo aggiornamento, Sez.~\ref{sec:aggiornamento})\\
\code{genera\_figure/} & 10 script (4 originali + 6 aggiunti)\\
\code{01\_genera\_dati\_vqe\_ideale.py} & script di generazione dati (VQE ideale)\\
\code{02\_scan36\_correlatori.py} & script di generazione dati (scan36, aggiunto)\\
\code{03\_correlatore\_da\_vqe.py} & script di generazione dati per \code{fig\_correlatore\_da\_vqe} (aggiunto)\\
\code{esegui\_tutto.py} & esegue tutto in un solo comando\\
\end{tabular}
\end{keyeq}

\subsection{I 9 moduli di \texttt{codice/}}

{\footnotesize
\begin{longtable}{@{}p{4.9cm}p{9.3cm}@{}}
\toprule
\textbf{File} & \textbf{Cosa fa} \\
\midrule
\endhead
\raggedright\code{dimer\_exact.py} & Hamiltoniana e diagonalizzazione esatta --- il ``mattone zero'' \\
\raggedright\code{vqe\_test2.py} & ansatz PMA-2q.3 e ottimizzazione VQE al punto di lavoro ``test 2'' \\
\raggedright\code{ansatz\_dimero.py} & libreria di varianti di ansatz (RBS, HA, PMA a diverse parametrizzazioni) \\
\raggedright\code{vqe\_dimer.py} & confronto sistematico HA vs PMA su tutto lo sweep in $B/J$ \\
\raggedright\code{trotter\_dimero.py} & evoluzione di Trotter per decomposizione esplicita in gate nativi \\
\raggedright\code{circuito\_correlazioni\_\allowbreak dimero.py} & circuito di Hadamard test per i correlatori dinamici \\
\raggedright\code{noise\_model\_dimero.py} & incluso solo per la costante \code{BASIS\_GATES} (vedi nota sotto) \\
\raggedright\code{stile.py} & stile grafico condiviso da tutti gli script di figura \\
\raggedright\code{validate\_ansatz\_params\_\allowbreak dimero.py} & controllo obbligatorio: preparazione VQE vs ampiezze esatte nei correlatori \\
\raggedright\code{validate\_circuito\_\allowbreak correlazioni\_dimero.py} & controllo obbligatorio: correlatori contro riferimento classico, convergenza di Trotter \\
\bottomrule
\end{longtable}
}

\begin{warnbox}{Una dipendenza reale verso il pacchetto rumoroso}
\code{validate\_ansatz\_params\_dimero.py} importa \code{BASIS\_GATES} da
\code{noise\_model\_dimero.py} --- un modulo concettualmente di Parte 2.
E' incluso qui \textbf{solo} per quella costante (la base di gate
$\{rz,sx,x,cx\}$ usata per transpilare), non per la sua funzionalità di
modellazione del rumore, mai usata in questo pacchetto.
\end{warnbox}

\section{Prerequisiti}

Stessi di \code{pacchetto\_dimero\_rumoroso\_riproducibile}:

\begin{keyeq}
\centering
\begin{tabular}{@{}ll@{}}
Python & 3.12.3\\
qiskit & 2.5.2\\
qiskit-aer & 0.17.2\\
numpy & 2.4.4\\
scipy & 1.17.1\\
matplotlib & 3.10.8\\
\end{tabular}
\end{keyeq}

\begin{okbox}{Nessuna installazione LaTeX richiesta}
Le figure sono generate direttamente in PDF/PNG da matplotlib.
\end{okbox}

\section{Il modo più semplice: un solo comando}

\begin{keyeq}
\centering
\code{cd pacchetto\_dimero\_riproducibile}\\
\code{python3 esegui\_tutto.py}
\end{keyeq}

Esegue in ordine: generazione dati (1 script), validazioni (2 script),
tutte le 4 figure. Si ferma al primo errore.

\begin{okbox}{Tempo stimato}
1--2 minuti. La fase più lenta è la figura di confronto ansatz (circa
30 secondi per 80 ottimizzazioni VQE in tutto).
\end{okbox}

Se tutto va a buon fine, l'ultima riga stampata è:
\begin{keyeq}
\centering
\code{TUTTO COMPLETATO. 4 figure in figure/, dati in dati/.}
\end{keyeq}

\section{Passo per passo, con la mappa esatta delle dipendenze}

Questa sezione non è solo un elenco di comandi: per ogni script indica
cosa prende in ingresso, cosa produce (file e contenuto), e chi
consuma quel prodotto più avanti nella catena. La buona notizia per
questo pacchetto: la catena è quasi piatta. Un solo script produce
dati; tre dei quattro script rimanenti (una validazione, tutte le
figure) \textbf{non hanno bisogno} di quel file --- calcolano da soli
tutto ciò che serve, in isolamento.

\subsection{Onda 1 --- nessuna dipendenza (possono girare in qualunque ordine, anche in parallelo)}

\begin{keyeq}[title={01\_genera\_dati\_vqe\_ideale.py}]
\textbf{Ingresso}: nessuno (parte da zero).\\
\textbf{Uscita}: \code{dati/ground\_state\_test2.npz} --- contiene
\code{b, J, D} (il punto di lavoro), \code{E0\_exact, psi0\_exact}
(benchmark esatto), \code{vqe\_params} (i 3 parametri ottimali
dell'ansatz PMA-2q.3), \code{fidelity, E\_vqe}.\\
\textbf{Chi lo consuma}: \code{validate\_ansatz\_params\_dimero.py}
(Onda 2, unico consumatore in questo pacchetto). Nota: è lo stesso file
prodotto in \code{pacchetto\_dimero\_rumoroso\_riproducibile} -- se lo hai
già generato là, puoi copiarlo qui invece di rieseguire lo script.\\
\textbf{Valore atteso}: energia $\approx-3.57321145$, fedeltà $\approx1.0000000000$.\\
\textbf{Comando} (dalla radice del pacchetto):\\
\code{python3 01\_genera\_dati\_vqe\_ideale.py}
\end{keyeq}

\begin{keyeq}[title={validate\_circuito\_correlazioni\_dimero.py}]
\textbf{Ingresso}: nessuno --- ricalcola da solo lo stato fondamentale
via \code{ground\_state()} di \code{circuito\_correlazioni\_dimero.py}.\\
\textbf{Uscita}: nessun file, solo un resoconto a schermo.\\
\textbf{Chi lo consuma}: nessuno (è un controllo terminale della catena).\\
\textbf{Comando} (dalla radice del pacchetto, nota il \code{cd}):\\
\code{cd codice}\\
\code{python3 validate\_circuito\_correlazioni\_dimero.py}\\
\code{cd ..}
\end{keyeq}

\begin{keyeq}[title={genera\_fig\_spettro\_esatto.py}]
\textbf{Ingresso}: nessuno --- diagonalizza da zero (nessun \code{.npz} richiesto).\\
\textbf{Uscita}: \code{figure/fig\_spettro\_dimero.pdf} (+\code{.png}).\\
\textbf{Chi la consuma}: nessuno (punto di arrivo).\\
\textbf{Comando} (dalla radice del pacchetto):\\
\code{python3 genera\_figure/genera\_fig\_spettro\_esatto.py}
\end{keyeq}

\begin{keyeq}[title={genera\_fig\_confronto\_ansatz.py}]
\textbf{Ingresso}: nessuno --- esegue da zero 80 ottimizzazioni VQE
(4 configurazioni $\times$ 20 punti in $b$).\\
\textbf{Uscita}: \code{figure/fig\_confronto\_ansatz\_HA\_PMA.pdf} (+\code{.png}).\\
\textbf{Chi la consuma}: nessuno (punto di arrivo).\\
\textbf{Comando} (dalla radice del pacchetto, $\sim$30 secondi):\\
\code{python3 genera\_figure/genera\_fig\_confronto\_ansatz.py}
\end{keyeq}

\begin{keyeq}[title={genera\_fig\_trotter\_convergenza.py}]
\textbf{Ingresso}: nessuno --- costruisce l'evoluzione via
\code{trotter\_dimero.py} da zero, a $b,J,D$ fissati nello script.\\
\textbf{Uscita}: \code{figure/fig\_trotter\_convergenza.pdf} (+\code{.png}).\\
\textbf{Chi la consuma}: nessuno (punto di arrivo).\\
\textbf{Comando} (dalla radice del pacchetto):\\
\code{python3 genera\_figure/genera\_fig\_trotter\_convergenza.py}
\end{keyeq}

\begin{keyeq}[title={genera\_fig\_correlatore\_convergenza.py}]
\textbf{Ingresso}: nessuno --- ricalcola lo stato fondamentale e il
riferimento classico da zero.\\
\textbf{Uscita}: \code{figure/fig\_correlatore\_convergenza.pdf} (+\code{.png}).\\
\textbf{Chi la consuma}: nessuno (punto di arrivo).\\
\textbf{Comando} (dalla radice del pacchetto):\\
\code{python3 genera\_figure/genera\_fig\_correlatore\_convergenza.py}
\end{keyeq}

\subsection{Onda 2 --- dipende dall'Onda 1}

\begin{keyeq}[title={validate\_ansatz\_params\_dimero.py}]
\textbf{Ingresso}: \code{ground\_state\_test2.npz} (copiato in
\code{codice/}) --- legge \code{vqe\_params} e \code{fidelity}.\\
\textbf{Uscita}: nessun file, solo un resoconto a schermo (include il
controllo 18 vs 19 CNOT).\\
\textbf{Chi lo consuma}: nessuno (controllo terminale).\\
\textbf{Comando} (dalla radice del pacchetto, dopo aver eseguito \code{01\_}):\\
\code{cp dati/ground\_state\_test2.npz codice/}\\
\code{cd codice}\\
\code{python3 validate\_ansatz\_params\_dimero.py}\\
\code{cd ..}
\end{keyeq}

\subsection{In sintesi: un solo ordine realmente obbligato}

\begin{okbox}{L'unica dipendenza vera di questo pacchetto}
\code{01\_genera\_dati\_vqe\_ideale.py} $\to$
\code{validate\_ansatz\_params\_dimero.py}. Tutto il resto (l'altra
validazione, tutte e 4 le figure) è indipendente e può girare prima,
dopo, o senza mai eseguire \code{01\_} --- \code{esegui\_tutto.py} li
esegue comunque tutti in un ordine fisso solo per chiarezza di lettura
del log, non perché serva.
\end{okbox}

\subsection{Tutti i comandi, in un unico blocco (nell'ordine di \texttt{esegui\_tutto.py})}

\begin{keyeq}
\centering
\begin{tabular}{@{}l@{}}
\code{cd pacchetto\_dimero\_riproducibile}\\
\code{python3 01\_genera\_dati\_vqe\_ideale.py}\\
\code{python3 02\_scan36\_correlatori.py}\\
\code{python3 03\_correlatore\_da\_vqe.py}\\
\code{cp dati/*.npz codice/}\\
\code{cd codice}\\
\code{python3 validate\_ansatz\_params\_dimero.py}\\
\code{python3 validate\_circuito\_correlazioni\_dimero.py}\\
\code{cd ..}\\
\code{python3 genera\_figure/genera\_fig\_confronto\_ansatz.py}\\
\code{python3 genera\_figure/genera\_fig\_correlatore\_convergenza.py}\\
\code{python3 genera\_figure/genera\_fig\_spettro\_esatto.py}\\
\code{python3 genera\_figure/genera\_fig\_trotter\_convergenza.py}\\
\code{python3 genera\_figure/genera\_fig\_gap\_magnetizzazione.py}\\
\code{python3 genera\_figure/genera\_fig\_vqe\_riparazione.py}\\
\code{python3 genera\_figure/genera\_fig\_trotter\_circuito\_vs\_matrice.py}\\
\code{python3 genera\_figure/genera\_fig\_scan36.py}\\
\code{python3 genera\_figure/genera\_fig\_correlatore\_da\_vqe.py}\\
\code{python3 genera\_figure/genera\_circuiti.py}\\
\end{tabular}
\end{keyeq}

\begin{okbox}{Risultati chiave dei due controlli}
\code{validate\_ansatz\_params\_dimero.py}: costo in gate, preparazione
esatta 18 CNOT contro preparazione VQE 19 CNOT (coincide esattamente
con quanto atteso).

\code{validate\_circuito\_correlazioni\_dimero.py}: convergenza
dell'errore di Trotter sul correlatore, rapporto che tende a 2
raddoppiando $N$ (Trotter al prim'ordine, errore $O(1/N)$).
\end{okbox}

Gli script di figura vanno eseguiti dalla \textbf{radice} del pacchetto,
non da dentro \code{genera\_figure/}.

\section{Aggiornamento: figure aggiunte}
\label{sec:aggiornamento}

Rispetto alla prima versione (4 figure, nessun diagramma di circuito),
aggiunte 15 figure per allineare il pacchetto al contenuto reale dei
quattro documenti narrativi (\texttt{tesi\_dimero\_slides}):

\begin{itemize}[leftmargin=1.4em,itemsep=3pt]
\item \textbf{Il gap esplicito}
(\code{genera\_fig\_gap\_magnetizzazione.py}): distanza
fondamentale-primo eccitato contro $B/J$, non solo lo spettro completo
--- minimo verificato $0.565$ con $D/J=0.2$, coincidente.
\item \textbf{La riparazione PMA}
(\code{genera\_fig\_vqe\_riparazione.py}): PMA-2q.3 (3 parametri)
raggiunge $\mathcal F>0.9998$ su tutto lo sweep, contro il crollo a
$\approx0.49$ di PMA a $K=1$ (Sez.~\ref{sec:pma-k1} sotto).
\item \textbf{Circuito vero contro matrice compatta}
(\code{genera\_fig\_trotter\_circuito\_vs\_matrice.py}): la prima
versione usava solo la matrice compatta (\code{U\_trotter\_mat}) come
unico metodo; qui entrambi i metodi sono calcolati e sovrapposti,
verificato che coincidono a precisione macchina ($\sim10^{-14}$), al
punto di lavoro $b/J=-0.18$, $D/J=1$ (stato $\ket{00}$, deliberatamente
asimmetrico --- vedi il documento per il perché).
\item \textbf{Scan36} (\code{02\_scan36\_correlatori.py} +
\code{genera\_fig\_scan36.py}): tutte le 36 combinazioni correlatore,
verificato 12/36 zeri a $t=0$, 0/36 zeri strutturali (nessuno resta
nullo per $\max_t$).
\item \textbf{Tutti i nove diagrammi di circuito}
(\code{genera\_circuiti.py}): assenti per scelta esplicita
dell'autore originale (vedi docstring di \code{ansatz\_dimero.py}).
Cinque riusano le chiamate esatte di un vecchio script del progetto
(\texttt{fig\_circuiti.py}, non incluso in questo pacchetto, letto come
riferimento); i restanti quattro (\code{circ\_pma\_base\_ramo2},
\code{circ\_rbs\_espanso}, \code{circ\_w\_espanso}, \code{circ\_cnotry})
non avevano uno script sorgente disponibile --- costruiti direttamente
dalle definizioni dei blocchi e dalle didascalie dei documenti.
\end{itemize}

\begin{okbox}{Risolto: la convenzione di fedeltà, non un bug --- uniformata al quadrato}
Il Documento 2 (\S6, tabella di confronto) riportava $\mathcal F=0.922857$
per l'ansatz PMA base (1 parametro) a $b/J=-0.18$, $D/J=1$, mentre la
ricostruzione qui arrivava a $0.851664$. Causa trovata: \code{vqe\_dimer.py}
calcolava $\mathcal F=\lvert\langle\psi_\text{esatto}\rvert\psi_\text{VQE}\rangle\rvert$
(modulo, senza quadrato) --- verificato $\sqrt{0.851664}=0.922856\ldots$,
coincide. Non era un errore di ricostruzione: era una convenzione di
fedeltà diversa da quella usata nel resto del progetto (dove si usa sempre
$\mathcal F=\lvert\langle\cdot\rvert\cdot\rangle\rvert^2$ --- inclusi
\code{vqe\_test2.py}, tutto il trimero, e \code{qiskit.quantum\_info.
state\_fidelity} di default). Uniformata qui al quadrato, la convenzione
più comune nel calcolo quantistico sperimentale/hardware: \code{vqe\_dimer.py}
corretto (riga 201), \code{dimero\_02\_vqe.tex} e \code{dimero\_03\_dinamica.tex}
aggiornati (5 punti: fig03, fig04, 3 tabelle di dettaglio), figure
rigenerate. La scelta dell'ansatz PMA-2q.3 al posto di PMA base non era in
discussione: elevare al quadrato una fedeltà $<1$ la rende solo più bassa,
mai più alta --- il caso per l'ansatz esteso è, se possibile, più netto con
la nuova convenzione, non più debole. \textbf{La figura
\code{fig14\_correlatore\_da\_vqe} è stata prodotta} con la convenzione
corretta (Sez.~\ref{sec:aggiornamento} sopra).

\textit{Correzione aggiuntiva, trovata allineando la griglia a quella di
\code{fig\_doc2.py} (lo script originale del Documento 2, letto come
riferimento):} con la griglia originale (20 punti, $[0,5]$)
\code{genera\_fig\_confronto\_ansatz.py} non campionava mai $B/J=2$
esattamente, dando un minimo apparente di $0.667$ invece del vero
$0.491178$. Corretta la griglia (33 punti, $[0,4]$, come l'originale) ---
ma questo ha esposto un problema reale: a $D=0$, $B/J=2$ il fondamentale è
\textbf{degenere}, e \code{vqe\_dimer.py} confrontava con un singolo
autovettore arbitrario da \code{eigh()}, dando fedeltà spurie vicine a
$0$ proprio dove fisicamente deve restare $1$. Corretto rendendo
\code{run\_vqe()} consapevole della degenerazione (proiezione sul
sottospazio fondamentale, stessa tecnica di \code{\_fidelity\_sottospazio}
in \code{fig\_doc2.py}, tol $10^{-9}$) --- verificato: $D=0$ ora dà
$\mathcal F=1$ ovunque per entrambi gli ansatz, $D/J=0.2$ dà il vero
minimo $0.491178$, coincidente con \code{fig\_doc2.py} e col documento.
\end{okbox}

Non ancora affrontate in questo aggiornamento: fig09 (ricco/piatto),
fig10/fig13 (ricostruzione spettrale), fig15/fig16 (scan36 con circuito
VQE reale, costoso computazionalmente).

\section{Un risultato genuino scoperto riproducendo la figura di confronto}
\label{sec:pma-k1}

\begin{warnbox}{PMA a $K=1$ non è sempre sufficiente}
Nella figura \code{fig\_confronto\_ansatz\_HA\_PMA}, la riga ``PMA
$D=0.2$'' mostra un calo netto di fedeltà (minimo $\approx0.49$, non
$1.0$) proprio nella zona di anticrossing ($B/J\approx2$) --- a
differenza di HA, che raggiunge sempre precisione macchina. Non è un
bug: con un solo parametro e uno stato iniziale scelto in modo rigido
(singoletto o $\ket{11}$ a seconda di $B/J$), l'ansatz PMA più semplice
non riesce a interpolare con continuità attraverso il crossover indotto
da $D\neq0$. E' precisamente il motivo per cui il resto del progetto usa
un ansatz PMA più espressivo (PMA-2q.3, 3 parametri, che rompe la
conservazione di $M$) invece di questa versione a un solo parametro.
\end{warnbox}

\section{Come sapere se un numero è quello giusto}

Ogni script stampa a schermo il valore appena calcolato \textbf{accanto}
al valore atteso.

\begin{warnbox}{Se un numero non coincide}
Qualcosa nel tuo ambiente locale sta producendo un risultato diverso:
fermati e confronta prima di proseguire.
\end{warnbox}

\section{Relazione con il pacchetto rumoroso}

\begin{okbox}{Stessa disciplina, stesso punto di partenza}
Questo pacchetto e \code{pacchetto\_dimero\_rumoroso\_riproducibile}
condividono lo script \code{01\_genera\_dati\_vqe\_ideale.py} e tre
moduli (\code{dimer\_exact.py}, \code{vqe\_test2.py},
\code{circuito\_correlazioni\_dimero.py}) --- non per caso: la Parte 2
(rumore) parte esattamente da dove finisce questa Parte 1, riusando lo
stesso punto di lavoro e la stessa preparazione VQE. Se hai già
eseguito il pacchetto rumoroso, \code{dati/ground\_state\_test2.npz} è
identico in entrambi.
\end{okbox}

\begin{okbox}{Verificato end-to-end, non solo assemblato}
L'intera pipeline è stata rieseguita da zero (cartelle \code{dati/} e
\code{figure/} svuotate, poi ogni script rilanciato in sequenza) dopo
ogni modifica sostanziale al codice --- ogni numero riportato in questa
guida è stato ricontrollato. Verificato di nuovo dopo l'uniformazione
della convenzione di fedeltà (Sez.~\ref{sec:aggiornamento}):
\code{exit code 0} su tutti gli script, \textbf{19 figure prodotte},
inclusa \code{fig\_correlatore\_da\_vqe} (ex \code{fig14}), non più
bloccata.
\end{okbox}

\end{document}
